\section*{Bab \ref{ch:b3:complete} --- Ruang Lengkap:
Baire, Ascoli, Stone--Weierstrass}

\begin{solution}{exo:b3:complete:1}
(a) Untuk $m \geq n$, selisih $f_n - f_m$ lenyap di luar sebuah selang
berpanjang $\frac1n$ dan terbatas oleh $1$: sehingga $\norm{f_n - f_m}_1
\leq \frac1n \to 0$: jadi Cauchy. Jika $f_n \to f$ dalam
$\norm\cdot_1$ dengan $f$ \omterm{def:b3:topology:continuity}{kontinu}: maka untuk $\alpha <
\frac12$ yang tetap, $\int_0^{\alpha}\abs f = \lim \int_0^\alpha\abs{f -
f_n + f_n} \leq \lim\bigl(\norm{f - f_n}_1 + 0\bigr) = 0$
(sebab $f_n \equiv 0$ di sana untuk $n$ besar), sehingga $f \equiv 0$ pada
$[0, \frac12)$ (menurut \omterm{def:b3:topology:continuity}{kekontinuannya}); demikian pula $f \equiv 1$ pada
$(\frac12, 1]$: dan tak ada fungsi \omterm{def:b3:topology:continuity}{kontinu} yang melakukan keduanya. Jadi ruangnya
tidak \omterm{def:b3:complete:complete}{lengkap} --- \omterm{thm:b3:complete:completion}{pelengkapannya} adalah $L^1$, yang dibangun pada
\cref{ch:b3:lp}.

(b) Misalkan $(x_n)$ Cauchy; pilihlah $n_1 < n_2 < \cdots$ dengan
$\norm{x_{n_{k+1}} - x_{n_k}} \leq 2^{-k}$. Maka deret
$\sum_k (x_{n_{k+1}} - x_{n_k})$ konvergen mutlak, sehingga
konvergen; jumlah parsialnya $x_{n_{k+1}} - x_{n_1}$, sehingga
$(x_{n_k})$ konvergen, dan barisan Cauchy yang mempunyai subbarisan
konvergen akan konvergen.
\end{solution}

\begin{solution}{exo:b3:complete:2}
(a) Andaikan $E$ \omterm{def:b3:complete:complete}{lengkap} dengan basis aljabar $(e_n)_{n\in\N}$
lalu misalkan $F_n = \operatorname{Vect}(e_1, \dots, e_n)$: yang tertutup
(sebab subruang berdimensi berhingga bersifat \omterm{def:b3:complete:complete}{lengkap}, sehingga tertutup ---
Tahun ke-2), dan berinterior kosong: sebab jika $B(x, r) \subseteq F_n$, ambillah
$v \notin F_n$; maka $x + \frac{r}{2\norm v}v \in B(x,r)
\setminus F_n$, mustahil. Padahal $E = \bigcup F_n$ (sebab setiap vektor berupa
kombinasi berhingga): dan itu bertentangan dengan Baire
(\cref{thm:b3:complete:baire}). Ruang $\R[X]$ mempunyai
basis terbilang $(X^n)$, sehingga tak ada norma yang membuatnya \omterm{def:b3:complete:complete}{lengkap}.

(b) Tetapkan $\delta > 0$ dan sebuah bola terbuka tak kosong $B$; kita cari di
$B$ sebuah titik $\Omega_\delta = \{x : \operatorname{osc}_xf <
\delta\}$, dengan $\operatorname{osc}_xf = \inf_{V \ni
x}\operatorname{diam} f(V)$. (Himpunan $\Omega_\delta$ terbuka: sebab jika
$\operatorname{diam}f(V) < \delta$ untuk suatu $V \ni x$ yang terbuka, maka setiap
$y \in V$ berosilasi $< \delta$.) Himpunan
\[
F_N = \bigl\{x : \abs{f_n(x) - f_m(x)} \leq \tfrac\delta3\
\ \forall\, n, m \geq N \bigr\}
\]
bersifat tertutup (sebagai irisan prapeta himpunan tertutup) dan
menyelimuti $\R$ (karena kekonvergenan titik demi titik membuat $(f_n(x))$ Cauchy).
Dengan menerapkan Baire di dalam $\bar B$ yang \omterm{def:b3:complete:complete}{lengkap}: ada $F_N \cap
\bar B$ yang memuat sebuah bola $B' = B(x_0, \rho)$. Dengan membiarkan $m \to
\infty$: $\abs{f_N - f} \leq \frac\delta3$ pada $B'$. Menurut
\omterm{def:b3:topology:continuity}{kekontinuan} $f_N$ di $x_0$, ciutkan menjadi $B'' \ni x_0$ tempat
$\abs{f_N - f_N(x_0)} \leq \frac\delta3$; maka untuk $x \in B''$,
\[
\abs{f(x) - f(x_0)} \leq \abs{f - f_N}(x) + \abs{f_N(x) -
f_N(x_0)} + \abs{f_N - f}(x_0) \leq \delta:
\]
sehingga $\operatorname{diam} f(B'') \leq 2\delta$, jadi $x_0 \in
\Omega_{3\delta} \cap B$. Dengan demikian setiap $\Omega_\delta$ terbuka dan
padat; lalu $\bigcap_k\Omega_{1/k}$ --- yakni himpunan titik \omterm{def:b3:topology:continuity}{kekontinuannya}
--- bersifat padat menurut Baire.
\end{solution}

\begin{solution}{exo:b3:complete:3}
(a) Pemetaan $f(x) = \frac12(x + \frac ax)$ memetakan $[\sqrt a, \infty)$ ke
dirinya sendiri (lewat AM--GM: $f(x) \geq \sqrt{x \cdot \frac ax} = \sqrt a$),
dan $f'(x) = \frac12(1 - \frac a{x^2}) \in [0, \frac12)$ di sana:
jadi kontraksi Lipschitz-$\frac12$ pada himpunan tertutup (yang \omterm{def:b3:complete:complete}{lengkap}).
Titik tetapnya: $x = f(x) \iff x^2 = a$: jadi $x^* = \sqrt a$. Lajunya
(\cref{thm:b3:complete:banach}): $d(x_n, \sqrt a) \leq
2^{-n+1}\,d(x_1, x_0)$. Untuk $a = 2$ dan $x_0 = 2$: $d(x_1, x_0) =
\frac12$, sehingga $n = 40$ menjamin $2^{-40} < 10^{-12}$. (Dalam
kenyataannya metode Newton konvergen secara kuadrat: segelintir
iterasi sudah cukup; taksiran kontraksinya pesimistis tetapi
gratis.)

(b) Pemetaan $f_m(x) = m + e\sin x$ bersifat Lipschitz-$e$ dengan $e < 1$ pada
$\R$ yang \omterm{def:b3:complete:complete}{lengkap}: jadi titik tetapnya tunggal, sebut $x(m)$. Untuk dua parameter:
\[
\abs{x(m) - x(m')} = \abs{f_m(x(m)) - f_{m'}(x(m'))}
\leq \abs{m - m'} + e\abs{x(m) - x(m')},
\]
sehingga $\abs{x(m) - x(m')} \leq \frac{\abs{m - m'}}{1 - e}$: jadi bahkan
Lipschitz terhadap $m$.
\end{solution}

\begin{solution}{exo:b3:complete:4}
(a) Misalkan $x^*$ titik tetap tunggal $f^p$. Maka
$f^p(f(x^*)) = f(f^p(x^*)) = f(x^*)$: sehingga $f(x^*)$ menjadi titik tetap
$f^p$, jadi $f(x^*) = x^*$. Dan titik tetap $f$ juga titik tetap
$f^p$: sehingga ketunggalannya berpindah. Untuk $T$: dengan induksi $\abs{T^pf(x)}
\leq \norm f_\infty x^p/p!$ (sebab setiap pengintegralan menambahkan faktor
$\frac xk$), sehingga $\norm{T^p} \leq a^p/p!$. Pemetaan $S(u) = g +
\lambda Tu$ memenuhi $S^p(u) - S^p(v) = \lambda^pT^p(u - v)$,
yang bernorma $\leq \abs\lambda^pa^p/p!\,\norm{u - v} \to 0$: jadi suatu
$S^p$ merupakan kontraksi, sehingga $S$ bertitik tetap tunggal: berarti
persamaan Volterra $u = g + \lambda Tu$ \omterm{def:b3:groups:derived}{terselesaikan} secara tunggal untuk
\emph{setiap} $\lambda$.

(b) Fungsi $\varphi(x) = d(x, f(x))$ \omterm{def:b3:topology:continuity}{kontinu} pada $K$ yang \omterm{def:b3:topology:compact}{kompak}:
sehingga ia mencapai minimumnya di suatu $x_0$. Jika $f(x_0) \neq x_0$,
maka $\varphi(f(x_0)) = d(f(x_0), f^2(x_0)) < d(x_0, f(x_0)) =
\min\varphi$: mustahil. Ketunggalannya: dua titik tetap $x \ne y$
akan memberikan $d(x,y) = d(f(x), f(y)) < d(x,y)$. Tanpa
kekompakan: $f(x) = x + \frac1x$ pada $[1, \infty)$ memenuhi,
untuk $x < y$, $f(y) - f(x) = (y - x)\bigl(1 - \frac1{xy}\bigr) <
y - x$, namun $f(x) > x$ selalu: jadi tak ada titik tetap --- sebab penurunan
jarak yang sejati lebih lemah daripada faktor kontraksi yang seragam.
\end{solution}

\begin{solution}{exo:b3:complete:5}
(a) Himpunan $A = \{g = h\}$ adalah prapeta diagonal
$\Delta_Y$ di bawah $x \mapsto (g(x), h(x))$ yang \omterm{def:b3:topology:continuity}{kontinu}; sedangkan $\Delta_Y$
tertutup karena $Y$ bersifat \omterm{def:b3:topology:hausdorff}{Hausdorff} (sebab untuk $y_1 \neq y_2$,
\omterm{def:b3:topology:topology}{persekitaran} terbuka yang saling lepas memberikan kotak terbuka di sekitar $(y_1,
y_2)$ yang saling lepas dengan diagonalnya, sehingga komplemen
$\Delta_Y$ terbuka): jadi $A$
tertutup, memuat \omterm{def:b3:topology:interior}{himpunan padat}, sehingga sama dengan $X$. Dipakai pada: ketunggalan di
\cref{thm:b3:complete:extension}, dan karenanya pada ketunggalan
\omterm{thm:b3:complete:completion}{pelengkapannya}.

(b) Karena $f$ merupakan isometri, ia \omterm{def:b3:topology:continuity}{kontinu} seragam: sehingga meluas menjadi
$F\colon X \to Y$ (\cref{thm:b3:complete:extension}), yang tetap
isometrik (sebab relasi $d(F(x), F(x')) = d(x, x')$ berlaku pada
himpunan pasangan yang padat dan kedua ruasnya \omterm{def:b3:topology:continuity}{kontinu}). Demikian pula
$f^{-1}\colon f(D) \to X$ meluas menjadi $G \colon Y \to X$. Komposisi
$G \circ F$ \omterm{def:b3:topology:continuity}{kontinu} dan menetapkan $D$ yang padat: jadi ia
$\mathrm{id}_X$ (menurut bagian (a)); secara simetris $F \circ G =
\mathrm{id}_Y$. Jadi $F$ merupakan bijeksi isometrik. Ketunggalan
\omterm{thm:b3:complete:completion}{pelengkapannya}: terapkan hal ini pada $D = X$ yang duduk secara padat di dua
\omterm{thm:b3:complete:completion}{pelengkapan}.
\end{solution}

\begin{solution}{exo:b3:complete:6}
Untuk $\{\sin(nx)\}$: terbatas titik demi titik oleh $1$; tetapi \emph{tidak}
\omterm{def:b3:complete:equicontinuous}{ekuikontinu}: sebab di $x = 0$ berlaku $\sin(n\cdot\frac{\pi}{2n}) = 1$
dengan $\frac\pi{2n} \to 0$, sehingga melanggar $\delta$ bersama mana pun untuk
$\varepsilon = \frac12$. Jadi tidak \omterm{def:b3:topology:compact}{kompak} relatif (menurut keperluan Ascoli,
\cref{thm:b3:complete:ascoli}).

Untuk $\{x^n\}$: terbatas; tetapi tidak \omterm{def:b3:complete:equicontinuous}{ekuikontinu} di $1$: sebab $1 - (1 -
\delta)^n \to 1$ ketika $n \to \infty$ untuk $\delta$ tetap. Jadi tidak
\omterm{def:b3:topology:compact}{kompak} relatif --- dan sejalan dengan itu, limit titik demi titiknya
tak \omterm{def:b3:topology:continuity}{kontinu}, sehingga tak ada subbarisan yang konvergen seragam.

Untuk $\{\norm f_\infty \leq 1, \norm{f'}_\infty \leq 1\}$: ketaksamaan
nilai rata-rata membuat keluarganya Lipschitz-$1$, sehingga
\omterm{def:b3:complete:equicontinuous}{ekuikontinu}; dan terbatas: jadi \omterm{def:b3:topology:compact}{kompak} relatif menurut Ascoli. (Tetapi tidak
\omterm{def:b3:topology:compact}{kompak}: sebab ia tidak tertutup --- limit seragamnya tak harus
$\mathcal C^1$; \omterm{def:b3:topology:interior}{penutupnya} adalah fungsi Lipschitz-$1$ yang
bernorma $\leq 1$.)

Untuk $\{\operatorname{Lip}(f) \leq 1, f(0) = 0\}$: \omterm{def:b3:complete:equicontinuous}{ekuikontinu};
terbatas titik demi titik (sebab $\abs{f(x)} \leq x \leq 1$); dan tertutup
terhadap limit seragam (sebab ketaksamaan Lipschitz dan nilainya di
$0$ beralih ke limitnya): jadi \omterm{def:b3:topology:compact}{kompak}.
\end{solution}

\begin{solution}{exo:b3:complete:7}
(a) Untuk $\norm f_\infty \leq 1$: $\abs{Tf(x)} \leq
\norm k_\infty$, dan
\[
\abs{Tf(x) - Tf(x')} \leq \int_0^1\abs{k(x,y) -
k(x',y)}\,\dd y \leq \omega_k\bigl(\abs{x - x'}\bigr),
\]
dengan $\omega_k$ berupa modulus \omterm{def:b3:topology:continuity}{kekontinuan} seragam $k$ pada
persegi \omterm{def:b3:topology:compact}{kompaknya} (menurut Heine): jadi peta bola satuannya
terbatas seragam dan \omterm{def:b3:complete:equicontinuous}{ekuikontinu}, sehingga \omterm{def:b3:topology:compact}{kompak} relatif
(menurut Ascoli). Lewat kelinearannya setiap himpunan terbatas berpeta \omterm{def:b3:topology:compact}{kompak}
relatif: jadi $T$ merupakan operator \omterm{def:b3:topology:compact}{kompak}.

(b) Pernyataan subbarisannya adalah definisi kekompakan
relatif yang diterapkan pada $\{Tf_n\}$. Seandainya $T$ bijektif dengan
invers yang \omterm{def:b3:topology:continuity}{kontinu}, maka $T(B(0,1)) \supseteq B(0,
\varepsilon)$ untuk suatu $\varepsilon > 0$ (sebab $T^{-1}$ \omterm{def:b3:topology:continuity}{kontinu}
di $0$); lalu bola tertutup $\bar B(0,\varepsilon)$, sebagai himpunan bagian
tertutup dari $\overline{T(B(0,1))}$ yang \omterm{def:b3:topology:compact}{kompak}, akan menjadi \omterm{def:b3:topology:compact}{kompak}
--- dan itu mustahil di $\mathcal C(\intcc01)$ yang berdimensi tak
berhingga menurut teorema Riesz (Tahun ke-2).
\end{solution}

\begin{solution}{exo:b3:complete:8}
(a) Dini: lihat \cref{lem:b3:complete:dini} --- yaitu argumen
penyelimutannya. (b) Salah: sebab gundukan berjalan $f_n(x) = \max(0, 1 -
\abs{nx - 1})$ menuju $0$ titik demi titik (sebab untuk $x > 0$, $f_n(x) =
0$ begitu $n > 2/x$; dan $f_n(0) = 0$) padahal $\norm{f_n}_\infty = 1$.
(c) Salah: sebab $f_n(x) = x^n$ turun ke $\mathbf 1_{\{1\}}$ yang tak
\omterm{def:b3:topology:continuity}{kontinu}; dan $\sup_{[0,1]}\abs{f_n - f} = \sup_{[0,1)}
x^n = 1$. (d) Salah: sebab $x^n \downarrow 0$ titik demi titik pada
$\intoo01$ yang tak \omterm{def:b3:topology:compact}{kompak}, dengan $\sup_{(0,1)}x^n = 1$. Jadi setiap
hipotesis Dini memang diperlukan.
\end{solution}

\begin{solution}{exo:b3:complete:9}
(a) Lewat kelinearannya $\int_0^1 fP = 0$ untuk setiap polinomial $P$.
Pilihlah $P_n \to f$ secara seragam
(\cref{cor:b3:complete:weierstrass}): maka $\int_0^1 f^2 =
\lim\int_0^1 fP_n = 0$, dan $f^2 \geq 0$ yang \omterm{def:b3:topology:continuity}{kontinu} dengan
integral nol pastilah lenyap secara identik.

(b) Pada $\intcc01$: polinomial dalam $x^2$ membentuk aljabar
yang memuat konstanta dan memisahkan titik (sebab $x \mapsto x^2$ injektif
pada $\intcc01$): jadi padat menurut Stone--Weierstrass. Pada $\intcc{-1}1$:
$x^2$ bernilai sama di $\pm x$, dan demikian pula setiap
polinomial dalam $x^2$: sehingga limit seragamnya berupa fungsi
genap. Seandainya fungsi genap $g_n$ konvergen seragam ke
identitasnya, maka $x = \lim g_n(x) = \lim g_n(-x) = -x$ untuk setiap
$x$: mustahil --- jadi tidak padat. Hipotesis pemisahannya gagal pada
pasangan $\{x, -x\}$.

(c) Tidak. Sebab untuk $f$ di dalam aljabar $\mathcal A$ yang dibangun oleh
konstanta dan $\eu^{\iu t}$ --- yakni kombinasi linear
$\eu^{\iu nt}$ dengan $n \geq 0$ --- berlaku $\Lambda(f) =
\int_0^{2\pi}f(t)\,\eu^{\iu t}\,\dd t = 0$ (sebab setiap
$\int_0^{2\pi}\eu^{\iu(n+1)t}\dd t = 0$). Fungsional $\Lambda$ bersifat
\omterm{def:b3:topology:continuity}{kontinu} terhadap $\norm\cdot_\infty$ (sebab $\abs{\Lambda(f)}\leq
2\pi\norm f_\infty$), sehingga $\Lambda$ lenyap pada $\bar{\mathcal
A}$; padahal $\Lambda(\eu^{-\iu t}) = 2\pi \neq 0$: jadi $\eu^{-\iu t}
\notin \bar{\mathcal A}$. (Stone--Weierstrass tak berlaku di sini:
sebab $\mathcal A$ tidak stabil terhadap konjugasi --- dan halangannya persis
yang akan disistematiskan teori fungsi holomorfik pada \cref{ch:b3:holomorphic}.)
\end{solution}

\begin{solution}{exo:b3:complete:10}
Himpunan $F_n = \bigcap_{f \in \mathcal F}\{x : \abs{f(x)} \leq n\}$ merupakan
irisan himpunan tertutup: jadi tertutup. Keterbatasan titik demi titiknya
memberikan $X = \bigcup_n F_n$. Lalu Baire (\cref{thm:b3:complete:baire})
menyediakan $n_0$ dengan $U = \mathring F_{n_0} \neq \varnothing$:
sehingga pada $U$ berlaku $\abs f \leq n_0$ untuk setiap $f \in \mathcal F$
serentak.
\end{solution}

\begin{solution}{exo:b3:complete:11}
(a) Barisan Cauchy terhadap $\norm\cdot_{\mathcal C^1}$ bersifat
Cauchy seragam bersama turunannya: sehingga $f_n \to f$
dan $f_n' \to g$ secara seragam, dengan $f, g$ \omterm{def:b3:topology:continuity}{kontinu}. Mengambil
limitnya (kekonvergenan seragam mengizinkannya di bawah
integral) pada $f_n(x) = f_n(0) + \int_0^xf_n'(t)\,\dd t$
memberikan $f(x) = f(0) + \int_0^xg$: jadi $f$ bersifat $\mathcal C^1$ dengan
$f' = g$, dan $\norm{f_n - f}_{\mathcal C^1} \to 0$.
Jadi \omterm{def:b3:complete:complete}{lengkap}.

(b) Fungsi $h(x) = \abs{x - \frac12}$ merupakan limit seragam
fungsi $\mathcal C^1$, misalnya $h_n(x) = \sqrt{(x -
\frac12)^2 + \frac1n}$ (sebab $\abs{h_n - h} \leq \frac1{\sqrt n}$
lewat batas kuantitas sekawan), padahal $h \notin E$: jadi
norma $\sup$ pada $E$ tidak \omterm{def:b3:complete:complete}{lengkap} --- \omterm{thm:b3:complete:completion}{pelengkapannya} adalah
$\mathcal C(\intcc01)$.

(c) Barisan $f_n(x) = \sin(2\pi nx)$ mempunyai $\norm{f_n}_\infty = 1$ dan
$\norm{f_n'}_\infty = 2\pi n \to \infty$: jadi tak ada $C$ yang berlaku.
(Secara konseptual: seandainya pendiferensialan terbatas terhadap norma
sup, kedua norma pada (a)--(b) akan setara, sehingga membuat
$(E, \norm\cdot_\infty)$ \omterm{def:b3:complete:complete}{lengkap} --- dan itu bertentangan dengan (b).
Ketakterbatasan inilah yang persis membuat fungsi \omterm{def:b3:topology:continuity}{kontinu} yang
umum dapat gagal diturunkan di mana pun,
\cref{thm:b3:complete:nowherediff}.)
\end{solution}

\begin{solution}{exo:b3:complete:12}
Tetapkan $\varepsilon > 0$. Setiap $F_N$ merupakan irisan atas $n
\geq N$ dari prapeta $\intcc{-\varepsilon}
\varepsilon$ yang tertutup di bawah $x \mapsto f(nx)$ yang \omterm{def:b3:topology:continuity}{kontinu}: jadi tertutup.
Hipotesisnya mengatakan bahwa setiap $x > 0$ terletak di suatu $F_N$. Menurut
Baire yang diterapkan di dalam $\intcc ab$ yang \omterm{def:b3:complete:complete}{lengkap} (untuk sembarang $0 < a <
b$), ada $F_N$ yang padat pada suatu subselang; dan karena tertutup, ia
memuat sebuah selang $\intcc{a'}{b'}$ dengan $0 < a' < b'$.
Lalu untuk setiap $n \geq \max(N, \frac{a'}{b' - a'})$, kedua
selang $\intcc{na'}{nb'}$ dan $\intcc{(n+1)a'}{(n+1)b'}$
bertindih (sebab $nb' \geq (n+1)a'$), sehingga
\[
\bigcup_{n \geq n_0}\intcc{na'}{nb'} \supseteq
\intco{n_0a'}{+\infty},
\]
dan setiap $t \geq n_0a'$ berbentuk $t = nx$ dengan $n \geq N$ dan $x \in
\intcc{a'}{b'} \subseteq F_N$: jadi $\abs{f(t)} \leq \varepsilon$.
Karena itu $\limsup_{t\to\infty}\abs f \leq \varepsilon$ untuk setiap
$\varepsilon$: sehingga $f \to 0$. Hipotesis penuhnya diperlukan
karena $F_N$ harus \emph{menyelimuti} sekumpulan $x$ sebesar selang
--- sebab dengan $x$ rasional saja gabungan $F_N$ bersifat
terbilang dan Baire tak memberikan apa pun; dan memang ada
contoh penyangkal yang \omterm{def:b3:topology:continuity}{kontinu}, yang lenyap sepanjang semua sinar
rasional tetapi tidak di tak hingga.
\end{solution}

\begin{solution}{pb:b3:complete:1}
\textbf{1.} Induksi pada $k$: jika $\norm{\varphi_n(t_k) - x_0}
\leq M(t_k - t_0) \leq MT \leq b$, maka titik $(t_k,
\varphi_n(t_k))$ terletak di $R$, sehingga kemiringan $f(t_k,
\varphi_n(t_k))$ terdefinisi dan bernorma $\leq M$; lalu untuk $t \in
[t_k, t_{k+1}]$ berlaku $\norm{\varphi_n(t) - x_0} \leq
\norm{\varphi_n(t_k) - x_0} + M(t - t_k) \leq M(t - t_0) \leq
MT \leq b$.

\textbf{2.} Setiap kepingan afinnya berkemiringan bernorma $\leq M$; dan
fungsi afin sepotong-sepotong yang kemiringannya terbatas oleh $M$ bersifat
Lipschitz-$M$ (rantaikan lewat titik kisinya).

\textbf{3.} Keluarga $(\varphi_n)$ terbatas titik demi titik
(sebab nilainya di $\bar B(x_0, b)$) dan \omterm{def:b3:complete:equicontinuous}{ekuikontinu} (dengan konstanta
Lipschitz bersama $M$): sehingga Ascoli
(\cref{thm:b3:complete:ascoli}) mengekstrak $\varphi_{n_j} \to
\varphi$ secara seragam pada $[t_0, t_0 + T]$. Batasnya beralih ke
limitnya: jadi $\varphi$ bersifat Lipschitz-$M$ dengan $\varphi(t_0) = x_0$.

\textbf{4.} Karena $R$ \omterm{def:b3:topology:compact}{kompak} dan $f$ \omterm{def:b3:topology:continuity}{kontinu}: maka ia \omterm{def:b3:topology:continuity}{kontinu}
seragam (menurut Heine, \cref{cor:b3:topology:heineborel}); misalkan
$\delta(\varepsilon)$ sebuah modulusnya. Untuk $t \in (t_k,
t_{k+1})$ yang bukan kisi: $\varphi_n'(t) = f(t_k, \varphi_n(t_k))$, sedangkan
kedua titik pemasukan nilai $f$ berbeda sebesar $\abs{t - t_k} \leq T/n$
pada waktunya dan $\norm{\varphi_n(t) - \varphi_n(t_k)} \leq MT/n$ pada
ruangnya. Jadi untuk $n \geq n_0$ dengan $(1 + M)T/n_0 < \delta(\varepsilon)$:
$\norm{\Delta_n(t)} \leq \varepsilon$.

\textbf{5.} Pada setiap $[t_k, t_{k+1}]$, $\varphi_n$ bersifat afin, sehingga
$\varphi_n(t_{k+1}) - \varphi_n(t_k) = \int_{t_k}^{t_{k+1}}
\varphi_n'(s)\dd s$, dengan turunannya berupa kemiringan yang konstan;
lalu menjumlahkan atas kepingannya (dan memotong yang terakhir di $t$):
$\varphi_n(t) = x_0 + \int_{t_0}^t\varphi_n'(s)\,\dd s$.
Menuliskan $\varphi_n' = f(s, \varphi_n(s)) + \Delta_n(s)$
(dengan integran \omterm{def:b3:topology:continuity}{kontinu} sepotong-sepotong, berlompatan berhingga banyak) memberikan
tampilan itu.

\textbf{6.} Diberikan $\varepsilon$: untuk $j$ besar berlaku
$\norm{\varphi_{n_j} - \varphi}_\infty < \delta(\varepsilon)$,
sehingga $\norm{f(s, \varphi_{n_j}(s)) - f(s, \varphi(s))} \leq
\varepsilon$ untuk setiap $s$: jadi integrannya konvergen
seragam, dan $\int_{t_0}^t f(s,\varphi_{n_j}(s))\dd s \to
\int_{t_0}^tf(s, \varphi(s))\dd s$ seragam terhadap $t$. Lagi pula
$\norm{\int_{t_0}^t\Delta_{n_j}} \leq T\sup\norm{\Delta_{n_j}}
\to 0$ (pertanyaan 4). Dengan mengambil limit pada identitas
pertanyaan 5: $\varphi(t) = x_0 + \int_{t_0}^tf(s,
\varphi(s))\,\dd s$.

\textbf{7.} Integran $s \mapsto f(s, \varphi(s))$ bersifat
\omterm{def:b3:topology:continuity}{kontinu}, sehingga ruas kanannya $\mathcal C^1$ terhadap $t$ dengan
turunan $f(t, \varphi(t))$: jadi $\varphi$ menyelesaikan masalah
Cauchy pada $[t_0, t_0 + T]$. Untuk separuh kirinya, tetapkan $g(t, x) =
-f(2t_0 - t, x)$, yang \omterm{def:b3:topology:continuity}{kontinu} pada persegi panjang cerminannya dengan
batas $M$ yang sama; lalu penyelesaian $\psi$ dari $y' = g(t,y)$,
$y(t_0) = x_0$ pada $[t_0, t_0 + T]$ melahirkan $\varphi(t) =
\psi(2t_0 - t)$ yang menyelesaikan persamaan aslinya pada $[t_0 - T,
t_0]$; dan kedua separuhnya merekat menjadi penyelesaian $\mathcal C^1$ (sebab kedua
turunan sepihaknya di $t_0$ sama dengan $f(t_0, x_0)$). ---
\emph{Inilah teorema Peano}: $f$ yang \omterm{def:b3:topology:continuity}{kontinu} mempunyai penyelesaian
lokal melalui setiap syarat awal.

\textbf{8.} \omterm{def:b3:topology:continuity}{Kekontinuannya} jelas. Sifat Lipschitz di dekat $0$ gagal:
sebab $\abs{f(h) - f(0)} = 2\sqrt h$, dan $2\sqrt h \leq L h$
salah untuk $h < 4/L^2$.

\textbf{9.} Untuk $t \leq c$: $x_c \equiv 0$ menyelesaikannya. Untuk $t \geq
c$: $x_c'(t) = 2(t - c) = 2\sqrt{(t-c)^2} = 2\sqrt{\abs{x_c}}$.
Di $t = c$ kedua turunan sepihaknya bernilai $0$: jadi $x_c$ bersifat
$\mathcal C^1$ dan menyelesaikannya secara global, dengan $x_c(0) = 0$ untuk setiap
$c \geq 0$ --- dan bersama $x \equiv 0$, terbentuklah kontinum
penyelesaian melalui titik asalnya.

\textbf{10.} Iterasi Picard menyusun $\Phi(x)(t) = x_0 +
\int_0^t f(x(s))\dd s$ dan menuntut $\norm{\Phi(x) - \Phi(y)}
\leq k\norm{x - y}$ dengan $k < 1$ pada bola yang sesuai --- dan itu
menyusul dari batas Lipschitz pada $f$, yang dipindahkan ke bawah
integralnya. Di sini $\sqrt\cdot$ tak punya batas Lipschitz di dekat $0$,
dan tak ada pilihan selang atau bola yang memperbaikinya. Dan memang tak ada
bukti ketunggalan yang bisa berhasil: sebab ketunggalannya \emph{salah}
(pertanyaan 9).

\textbf{11.} Di $c_0$, vektor satuan $e_n = (0, \dots, 0, 1,
0, \dots)$ memenuhi $\norm{e_n - e_m}_\infty = 1$ untuk $n \neq
m$: sehingga tak ada subbarisan yang Cauchy, jadi bola satuan tertutupnya tidak
\omterm{def:b3:topology:compact}{kompak}. Langkah yang runtuh adalah ekstraksinya (pertanyaan 3):
sebab Ascoli untuk $\mathcal C([t_0, t_0+T], E)$ menuntut nilainya
hidup di ruang tempat himpunan terbatas bersifat \omterm{def:b3:topology:compact}{kompak} relatif ---
benar di $\R^d$ (lewat Bolzano--Weierstrass), salah di $c_0$; sehingga ekstraksi
titik demi titik tak lagi tersedia (dan memang contoh Dieudonné
sama sekali tak punya penyelesaian lokal).

\textbf{12.} \emph{Peano}: jika $f$ \omterm{def:b3:topology:continuity}{kontinu} pada \omterm{def:b3:topology:topology}{persekitaran}
$(t_0, x_0)$ di $\R\times\R^d$, $\Rightarrow$ maka ada penyelesaian
$\mathcal C^1$ bagi $x' = f(t,x)$, $x(t_0) = x_0$, pada
suatu $[t_0 - T, t_0 + T]$. \emph{Cauchy--Lipschitz}
(\cref{ch:b3:ode}): jika di samping itu $f$ bersifat Lipschitz lokal terhadap
variabel $x$, maka penyelesaiannya \emph{tunggal} (sebab dua penyelesaian bersesuaian pada
selang bersamanya) --- yakni keberadaan sekaligus ketunggalan. Pasangan
$(x' = 2\sqrt{\abs x},\ x(0) = 0)$ memisahkan kedua teorema itu.

\textbf{13.} Untuk $\omega(r) = Lr$: $\int_0^1\frac{\dd r}{Lr} =
+\infty$: jadi memenuhi. Untuk $\omega(r) = r\log\frac1r$ (di dekat $0$):
$\int\frac{\dd r}{r\log(1/r)} = \bigl[-\log\log\frac1r\bigr]
\to +\infty$ ketika $r \to 0$: jadi memenuhi --- padahal $\omega(r)/r =
\log\frac1r \to \infty$: jadi tidak Lipschitz. Untuk $\omega(r) = 2\sqrt
r$: $\int_0^1\frac{\dd r}{2\sqrt r} = \bigl[\sqrt r\bigr]_0^1
= 1 < \infty$: jadi gagal.

\textbf{14.} Kurangkan kedua bentuk integralnya:
\[
x_1(t) - x_2(t) = x_1(s) - x_2(s) + \int_s^t\bigl(f(v,
x_1(v)) - f(v, x_2(v))\bigr)\dd v,
\]
lalu ambil normanya dan pakai modulus Osgoodnya: sehingga $\delta(t) \leq
\delta(s) + \int_s^t\omega(\delta(v))\,\dd v$.

\textbf{15.} Bilangan $\tau$ terdefinisi dengan baik (sebab $\delta(t_0) = 0$) dengan
$\delta(\tau) = 0$ (menurut \omterm{def:b3:topology:continuity}{kekontinuannya}) dan $\delta > 0$ pada
$\intoc\tau{t_1}$. Tetapkan $s \in \intoo\tau{t_1}$. Maka $u$ bersifat
$\mathcal C^1$, $u(s) = \delta(s) > 0$, $u \geq \delta$ pada
$[s, t_1]$ (pertanyaan 14), dan $u' = \omega(\delta) \leq
\omega(u)$ (sebab $\omega$ tak turun dan $u > 0$). Bagilah lalu
integralkan:
\[
\int_{\delta(s)}^{u(t_1)}\frac{\dd r}{\omega(r)}
= \int_s^{t_1}\frac{u'(v)}{\omega(u(v))}\,\dd v \leq t_1 - s
\leq t_1 - \tau .
\]
Meski $u$ bergantung pada $s$, batas $u(t_1) \geq
\delta(t_1)$ berlaku untuk setiap $s$, sehingga ruas kirinya sekurang-kurangnya
$\int_{\delta(s)}^{\delta(t_1)}\frac{\dd r}{\omega(r)}$,
yang menuju $+\infty$ ketika $s \downarrow \tau$ (sebab lalu
$\delta(s) \to \delta(\tau) = 0$, dan integralnya divergen
di $0$): jadi ruas kanannya yang terbatas terbantahkan. Karena itu
$\delta \equiv 0$: yakni ketunggalannya.

\textbf{16.} Lipschitz adalah kasus $\omega = Lr$: sehingga ketunggalannya diperoleh kembali.
Untuk $x' = x\log\frac1{\abs x}$: ruas kanannya memenuhi
modulus Osgood $\omega(r) = r\log\frac1r$ di dekat $0$ (lewat ketaksamaan
nilai rata-rata pada $x \mapsto x\log\frac1x$, yang turunannya
$\log\frac1x - 1$ tak terbatas --- jadi Lipschitz gagal, Osgood
berlaku): sehingga penyelesaiannya tunggal; catat bahwa $x \equiv 0$ salah satunya, sehingga tak ada
penyelesaian lain yang dapat menyentuh $0$. Untuk $\omega(r) = 2\sqrt r$:
$\int_0^{h^2}\frac{\dd r}{2\sqrt r} = h$ --- jadi ``anggaran
Osgood'' untuk mendaki dari $0$ ke ketinggian $h^2$ tepat sebesar
waktu $h$, dan memang $x_c(c + h) = h^2$: penyelesaiannya memakai
waktu $h$ untuk melakukan persis apa yang diizinkan integral yang
konvergen itu. Divergennya integral berarti kemustahilan
meninggalkan $0$ dalam waktu berhingga; sedangkan konvergennya adalah jalan lolosnya.

\textbf{17.} Fungsi $G(t) = \int_{t_0}^t(Le(s) + \eta(s))\dd s$ bersifat
$\mathcal C^1$ dengan $G' = Le + \eta \leq LG + \sup\eta$
(menurut hipotesis $e \leq G$). Maka $\bigl(\eu^{-Lt}(G +
\tfrac{\sup\eta}L)\bigr)' = \eu^{-Lt}(G' - LG - \sup\eta)
\leq 0$: sehingga kurungnya menurun, jadi $G(t) + \frac{\sup\eta}L
\leq \eu^{L(t - t_0)}\bigl(G(t_0) + \frac{\sup\eta}L\bigr) =
\eu^{L(t-t_0)}\frac{\sup\eta}L$, yakni $e(t) \leq G(t) \leq
\frac{\sup\eta}L(\eu^{L(t-t_0)} - 1)$.

\textbf{18.} Cacat kuantitatifnya: pada $(t_k, t_{k+1})$ berlaku
$\Delta_n(t) = f(t_k, \varphi_n(t_k)) - f(t, \varphi_n(t))$
dengan $\abs{t - t_k} \leq T/n$ dan $\norm{\varphi_n(t) -
\varphi_n(t_k)} \leq MT/n$, sehingga $\norm{\Delta_n} \leq L'T/n +
LMT/n$. Dengan mengurangkan identitas integral $\varphi_n$
(pertanyaan 5) dan $\varphi$ lalu memakai batas Lipschitznya:
\[
\norm{\varphi_n(t) - \varphi(t)} \leq
\int_{t_0}^t L\,\norm{\varphi_n - \varphi}(s)\,\dd s +
\int_{t_0}^t\norm{\Delta_n(s)}\,\dd s,
\]
lalu pertanyaan 17 dengan $\eta = \norm{\Delta_n}$ memberikan
batas $O(1/n)$ yang dinyatakan itu, dengan $\varphi$ yang tunggal menurut
Cauchy--Lipschitz (atau Osgood). Bahkan tanpa laju: setiap
subbarisan dari $(\varphi_n)$ yang terbatas dan Lipschitz secara seragam
mempunyai sub-subbarisan yang konvergen (lewat Ascoli + Bagian II) ke
\emph{suatu} penyelesaian, yang oleh ketunggalannya terpaksa menjadi $\varphi$:
dan barisan yang setiap subbarisannya mempunyai sub-subbarisan
berlimit sama pastilah konvergen.

\textbf{19.} Dari $\varphi_n(t_k) = 0$: kemiringan $2\sqrt0 =
0$ memberikan $\varphi_n(t_{k+1}) = 0$; lalu dengan induksi $\varphi_n
\equiv 0$, sehingga konvergen ke penyelesaian nolnya. Dari
$\varepsilon > 0$: pada $[\varepsilon', \infty)$ dengan
$\varepsilon' < \varepsilon$ fungsi $2\sqrt x$ bersifat
Lipschitz, sehingga pertanyaan 18 berlaku dan Euler konvergen ke
penyelesaian tunggal melalui $(0, \varepsilon)$, yakni $x(t) =
(t + \sqrt\varepsilon)^2$ (periksa: $x' = 2(t + \sqrt
\varepsilon) = 2\sqrt x$). Ketika $\varepsilon \to 0$, $(t +
\sqrt\varepsilon)^2 \to t^2$ seragam pada $[0,1]$: jadi limit
gandanya mendarat pada $x_0(t) = t^2$, bukan pada $0$. Usikan
sekecil apa pun pada data awalnya mengalihkan skemanya
dari satu penyelesaian ke penyelesaian lain: jadi ketaktunggalannya terbaca sebagai
ketakstabilan numerik.

\textbf{20.} Ketakkosongannya: lihat Bagian II. Setiap penyelesaian memenuhi
$\norm{x'} = \norm{f(t, x)} \leq M$: sehingga $\mathcal S$ bersifat
Lipschitz-$M$ secara seragam dan terbatas seragam (sebab nilainya di
$\bar B(x_0, b)$). Ketertutupannya: jika $x_n \in \mathcal S \to x$
secara seragam, ambillah limit pada $x_n(t) = x_0 +
\int_{t_0}^tf(s, x_n(s))\dd s$ (sebab integrannya konvergen
seragam menurut \omterm{def:b3:topology:continuity}{kekontinuan} seragam $f$ pada $R$ yang \omterm{def:b3:topology:compact}{kompak}):
maka $x \in \mathcal S$. Lewat Ascoli: $\mathcal S$ merupakan himpunan bagian
$\mathcal C([t_0, t_0+T], \R^d)$ yang tertutup, terbatas, dan \omterm{def:b3:complete:equicontinuous}{ekuikontinu}: jadi \omterm{def:b3:topology:compact}{kompak}.

\textbf{21.} Setiap penyelesaian bersifat tak turun (sebab $x' = 2\sqrt{\abs
x} \geq 0$) dengan $x(0) = 0$, sehingga $x \geq 0$. Misalkan $c =
\sup\{t \in [0,1] : x(t) = 0\}$ (yang mungkin $c = \infty$ bila
$x \equiv 0$, dan dalam hal itu $x = x_\infty$). Untuk $t > c$: $x
> 0$ (menurut kemonotonannya dan definisi $c$), dan di sana $(\sqrt
x)' = \frac{x'}{2\sqrt x} = 1$, sehingga $\sqrt{x(t)} = t - c$
(menurut \omterm{def:b3:topology:continuity}{kekontinuannya} di $c$): jadi $x = x_c$. \omterm{def:b3:topology:continuity}{Kekontinuan} $c \mapsto
x_c$: untuk $c, c' \in [0, 1]$ berlaku $\sup_{[0,1]}\abs{(t - c)_+^2 -
(t - c')_+^2} \leq 2\abs{c - c'}$ (sebab pemetaan $c \mapsto
(t-c)_+^2$ bersifat Lipschitz-$2$ seragam terhadap $t \in [0,1]$), dan
karena $x_c \equiv 0$ pada $[0,1]$ untuk setiap $c \geq 1$, keluarganya
menyusut menjadi $\mathcal S = \{x_c : c \in [0,1]\}$ (dengan
$x_1 = 0 = x_\infty$). Jadi $\mathcal S$ merupakan peta
$[0,1]$ yang \omterm{def:b3:topology:compact}{kompak} dan \omterm{def:b3:topology:connected}{terhubung} di bawah \omterm{def:b3:topology:continuity}{pemetaan kontinu} $c
\mapsto x_c$: sehingga \omterm{def:b3:topology:compact}{kompak} dan \omterm{def:b3:topology:connected}{terhubung}. Corong penyelesaiannya
berupa ruas garis \omterm{def:b3:topology:continuity}{kontinu} yang membentang dari $t^2$ (lolos
seketika) turun ke $0$ (istirahat abadi).

\textbf{22.} Pemasukan nilai $\operatorname{ev}_t\colon
\mathcal C([0,1]) \to \R$, $x \mapsto x(t)$, bersifat \omterm{def:b3:topology:continuity}{kontinu}
(sebab $\abs{x(t) - y(t)} \leq \norm{x - y}_\infty$), sehingga $\mathcal
S(t) = \operatorname{ev}_t(\mathcal S)$ merupakan peta \omterm{def:b3:topology:continuity}{kontinu}
sebuah \omterm{def:b3:topology:compact}{kompak}: jadi \omterm{def:b3:topology:compact}{kompak} --- dan peta himpunan \omterm{def:b3:topology:connected}{terhubung}: jadi \omterm{def:b3:topology:connected}{terhubung}.
Untuk contohnya: $x_c(t) = (t - c)_+^2$ menyapu, ketika $c$ menjelajahi
$[0, 1]$, semua nilai dari $t^2$ (di $c = 0$) turun ke
$0$ (di $c \geq t$), secara \omterm{def:b3:topology:continuity}{kontinu}: jadi $\mathcal S(t) =
\intcc0{t^2}$. Pada setiap saat, penampang corongnya berupa
ruas garis penuh: jadi di antara istirahat dan pelolosan maksimal, setiap
kompromi diwujudkan oleh sebuah penyelesaian yang sesungguhnya.

\textbf{23.} Dengan mengurangkan bentuk integral $x(t) = x_0 +
\int_{t_0}^tf(s, x(s))\dd s$ dan padanannya untuk $y$, lalu
menetapkan $e(t) = \norm{x(t) - y(t)}$ dan $d = \norm{x_0 - y_0}$:
\[
e(t) \leq d + \int_{t_0}^tL\,e(s)\dd s = G(t).
\]
Maka $G(t_0) = d$ dan $G' = Le \leq LG$, sehingga $\bigl(\eu^{-L(t
- t_0)}G\bigr)' \leq 0$ dan $G(t) \leq d\,\eu^{L(t - t_0)}$;
karena itu $e(t) \leq d\,\eu^{L(t-t_0)}$. Ketajamannya: untuk $x' =
Lx$, penyelesaian melalui $x_0$ dan $y_0$ adalah $x_0\eu^{L(t -
t_0)}$ dan $y_0\eu^{L(t-t_0)}$, yang jaraknya tepat
$d\,\eu^{L(t-t_0)}$. Dengan $d = 0$ diperoleh $e \equiv 0$: yakni
ketunggalan Cauchy--Lipschitz, yang diturunkan ulang dalam dua baris. Dan
untuk $t \in [t_0, t_0 + T]$ yang tetap, $\norm{x(t; x_0) - x(t;
y_0)} \leq \eu^{LT}\norm{x_0 - y_0}$: jadi alirannya bersifat Lipschitz
terhadap syarat awalnya --- yakni kebergantungan deterministik, dengan
harga eksponensial yang terkendali.

\textbf{24.} Pada $\intoc01$, $\Omega$ terdefinisi dengan baik (sebab
integralnya konvergen di $0$ menurut hipotesis) dan bersifat $\mathcal C^1$ dengan
$\Omega' = \frac1\omega > 0$: jadi bijeksi naik pada
$\intoc0{\Omega(1)}$, dengan $\Omega(x) \to 0$ ketika $x \to 0^+$.
Inversnya $g \colon \intoc0{\Omega(1)} \to \intoc01$ bersifat
$\mathcal C^1$ dengan
\[
g'(t) = \frac1{\Omega'(g(t))} = \omega\bigl(g(t)\bigr) > 0,
\qquad g(t) \xrightarrow[t \to 0^+]{} 0 .
\]
Perluaslah $g$ dengan $0$ pada $t \leq 0$: \omterm{def:b3:topology:continuity}{kekontinuannya} jelas, dan di
$t = 0$, untuk $t > 0$,
\[
\frac{g(t)}t = \frac1t\int_0^tg'(s)\dd s
= \frac1t\int_0^t\omega\bigl(g(s)\bigr)\dd s
\leq \omega\bigl(g(t)\bigr) \xrightarrow[t\to0^+]{} 0
\]
(sebab $\omega$ tak turun, $g$ naik, dan $\omega(0^+) = 0$):
sehingga $g'(0) = 0 = f(g(0))$, dan $g' = f(g)$ berlaku pada kedua sisi
$0$. Jadi $x \equiv 0$ dan $g$ merupakan dua penyelesaian $\mathcal
C^1$ berbeda yang melalui $(0, 0)$: jadi ketika $\int_0\frac{\dd r}
{\omega(r)}$ konvergen, ketunggalannya gagal --- sehingga divergensi pada
pertanyaan 15 tepat merupakan perbatasannya. Untuk $\omega(r) = 2\sqrt
r$: $\Omega(x) = \int_0^x\frac{\dd r}{2\sqrt r} = \sqrt x$,
$g(t) = t^2$, dan translasi waktu memberikan seluruh keluarga
$x_c$ pada Bagian III.

\textbf{25.} Dengan langkah $h = \frac1n$: $\varphi_n(t_{k+1}) =
\varphi_n(t_k) + h\,\varphi_n(t_k) = (1 + h)\varphi_n(t_k)$,
sehingga $\varphi_n(1) = (1 + \frac1n)^n$ setelah $n$ langkah.
Penjabarannya:
\[
n\log\Bigl(1 + \frac1n\Bigr)
= n\Bigl(\frac1n - \frac1{2n^2} + O\Bigl(\frac1{n^3}\Bigr)
\Bigr) = 1 - \frac1{2n} + O\Bigl(\frac1{n^2}\Bigr),
\]
lalu dengan mengeksponenkannya, $(1 + \frac1n)^n = \eu\,\eu^{-1/(2n) +
O(n^{-2})} = \eu\bigl(1 - \frac1{2n} + O(n^{-2})\bigr)$. Jadi
galat di $t = 1$ adalah $\eu - \varphi_n(1) =
\frac{\eu}{2n} + O(n^{-2})$: yakni $O(\frac1n)$ pada pertanyaan 18,
di sini beserta konstanta eksaknya $\frac\eu2$. Secara numerik, untuk $n =
10$: $1.1^{10} = 2.5937424601$ (sebab $1.1^2 = 1.21$, $1.1^4 =
1.4641$, $1.1^8 = 2.14358881$, dikalikan $1.21$), dan $\eu -
2.59374 \approx 0.12454$, berbanding ramalan asimtotiknya
$\frac{\eu}{20} \approx 0.13591$: jadi sesuai dalam batas
koreksi $O(n^{-2})$, yang suku utamanya di sini menurunkan
ramalannya menuju nilai yang teramati.

\end{solution}
